%% ============================================================================
%% sections/05-methodology.tex -- Evaluation Methodology
%% Class-independent section file: it is \input by main-acm.tex and main-ieee.tex.
%% Rules (see README-build.md): no class-specific commands; cite with \cite only;
%% figure sizes come from the float budget and table bodies from tables/, never from this file.
%% Every paragraph and every caption is one numbered unit that starts with a
%% marker line "%% P:<id>".  Ids are frozen: never renumber, never reuse.
%% The same ids appear in paper.zh-CN.md; check_align.py compares the two.
%% Word budget of this section: 700 words of running text.
%% ============================================================================

\section{Evaluation Methodology}\label{sec:method}

%% P:fig8-cap
\FigSlot{fig:flow}%
  {Flow diagram with two columns of boxes. Left, the inputs: golden model plus tests, synthetic tensors, RTL plus configuration, flow plus PDK setup, and workload inventory. Right, the forms of evidence: RTL simulation, synthesis slice, post-route EDA estimate, analytical derivation and, in a dotted box, model results. Solid lines join the golden model, the synthetic tensors and the RTL to RTL simulation; the RTL and the flow setup to synthesis slice and to post-route EDA estimate; and the workload inventory to analytical derivation. A dotted line runs from the workload inventory through a dotted box marked model, to be defined, to the model results}%
  {drawn; file figures/fig-flow.pdf; the modeling step is an open slot}%
  {Planned evaluation flow. Left: inputs. Right: the form of evidence that labels each number. Workload inventories are derived analytically and do not pass through a model. Dotted parts are open: the model has not been chosen, and only one of its inputs is drawn. The numerical front end that runs the synthetic tensors is not drawn.}

\subsection{Configurations and Evidence Types}\label{sec:configs}

%% P:tab4-cap
\begin{table}[t]
\centering\SlotName{tab:configs}%
\caption{Configurations and forms of evidence.}\label{tab:configs}
%% tables/tab-configs.tex -- body of the table `tab:configs'.
%% The float, the caption and the label are in sections/05-methodology.tex; this file sets the type size and the
%% column separation itself, so that it alone decides how wide and how high the table is.
%% No blank line in this file: check_align.py reads it as part of the unit that inputs it.
\TableSize\setlength{\tabcolsep}{3pt}%
\begin{tabular}{@{}ll@{}}
\toprule
Configuration & Definition \\
\midrule
\cfgAsBuilt & RTL as it stands; behavioral dual-port \Lii in simulation \\
\cfgCorrected & plus a versioned list of fixes; main experiments \\
\cfgProjected & plus more PEs, network and shared storage; projections \\
(listed apart) & target-design points not yet in RTL \\
\midrule
\multicolumn{2}{@{}p{\PaperColW}@{}}{Forms of evidence: \evRTL, \evModel, \evAnalytic, \evSlice, \evPostRoute.} \\
\bottomrule
\end{tabular}

\end{table}

%% P:sec5.1-p1
\ptag{sec5.1-p1}%
Every number carries two labels: one of three configurations \TBD{N}{neutral names for the three configurations} and a form of evidence (\tabref{tab:configs}, \figref{fig:flow}).
The fixes of \cfgCorrected include error reporting to the command processor, the contract of the single-port \Lii macros and scale write-back \TBD{D7}{final list of fixes, pending the merge plan}.
A physical result is never called measured.

\subsection{Metrics, Operating Regimes and \texorpdfstring{$K$}{K}}\label{sec:metrics}

%% P:sec5.2-p1
\ptag{sec5.2-p1}%
We report latency and throughput at three scopes: the end-to-end period of the complete chain, the steady-state throughput of useful results, and the local period of the \handoff and consumption stage.
A local gain does not imply an end-to-end gain, and end-to-end parity does not show equal \handoff cost.
The other metrics are area and timing, SRAM transactions, numerical error against the common error budget, and energy as action counts times sourced unit costs plus leakage, left open where a cost has no source.
Two operating regimes are covered: \Lii-resident, with operands already in local storage, and supply-limited, where the PE port bounds their delivery \TBD{D4}{definition of the supply-limited regime at the planned port width}.
The total $K$ of a GEMM is kept apart from the $K_c$ of one tile, and placements are compared over the range of $K$ that is legal for all of them \TBD{D7}{legal range of K under the target operand loader}.

\subsection{Numerical Protocol}\label{sec:numerics}

%% P:sec5.3-p1
\ptag{sec5.3-p1}%
Numerical order is frozen before performance is compared (\secref{sec:placements}), down to the final consumption point.
A bit-accurate reference (golden) model defines the hardware semantics; all placements get the same tensors and seeds.
The inputs are synthetic: we cannot run the models of \secref{sec:workload-method} and do not use their weights.
Shapes come from the workload's matrix products; activations and cache entries follow the quantization rules of the reference implementation, which gives none for weights \TBD{E3}{value distributions and weight quantization rule of the synthetic tensors}.
Quality means numerical error against a reference-precision result, not a task metric.

\subsection{Workload: What Is Counted and How}\label{sec:workload-method}

%% P:tab5-cap
\begin{table}[t]
\centering\SlotName{tab:mac-share}%
\caption{Share of MACs per decoded token by operand data type of the reference implementations (\evAnalytic).}\label{tab:mac-share}
%% tables/tab-mac-share.tex -- body of the table `tab:mac-share'.
%% The float, the caption and the label are in sections/05-methodology.tex; this file sets the type size and the
%% column separation itself, so that it alone decides how wide and how high the table is.
%% No blank line in this file: check_align.py reads it as part of the unit that inputs it.
\TableSize\setlength{\tabcolsep}{2.5pt}%
\begin{tabular}{@{}lr@{\,/\,}lr@{\,/\,}lr@{\,/\,}l@{}}
\toprule
Operands; share at 8K / 1M (\%) & \multicolumn{2}{c}{V4-Pro} & \multicolumn{2}{c}{V4-Flash} & \multicolumn{2}{c@{}}{V4.1-Flash} \\
\midrule
4 $\times$ 4 bit (index scoring) & 0.9 & 42.5 & 2.4 & 64.3 & 1.2 & 38.3 \\
4 $\times$ 8 bit (routed experts)\textsuperscript{a} & 44.2 & 15.9 & 44.0 & 9.3 & 47.3 & 29.5 \\
8 $\times$ 8 bit (projections, shared expert) & 34.0 & 12.3 & 29.9 & 6.3 & 30.6 & 19.1 \\
BF16 $\times$ cache (attention core) & 9.7 & 25.3 & 7.8 & 16.8 & 9.0 & 5.6 \\
BF16 $\times$ BF16 (output projection)\textsuperscript{b} & 7.5 & 2.7 & 9.8 & 2.1 & 7.5 & 4.7 \\
FP32 (remaining matrices) & 3.6 & 1.3 & 6.2 & 1.3 & 4.4 & 2.8 \\
\bottomrule
\multicolumn{7}{@{}l@{}}{\textsuperscript{a}\,4-bit weights. \textsuperscript{b}\,Stored 8-bit. V4.1-Flash at 1M: its report's count.} \\
\end{tabular}

\end{table}

%% P:sec5.4-p1
\ptag{sec5.4-p1}%
The workload is the decode step of the three models of \secref{sec:bg-workload} \cite{DeepSeekAI2026V4ProHF,DeepSeekAI2026V41FlashHF} \citeTBD{DeepSeek-V4-Flash model repository}; V4-Flash is counted with the V4-Pro code under its own configuration.
After prefill the reference implementation feeds one token per step, so counts are per sequence and per new token at context length $L$.
MACs are counted only inside matrix products, other operations separately.

%% P:sec5.4-p2
\ptag{sec5.4-p2}%
The counts are an \evAnalytic; no model is run.
Each matrix product contributes rows times inner-product length times layers, and other operations follow the loop bounds in the source.
Operand types follow the reference code (\tabref{tab:mac-share}); index scoring is counted as 4-bit by 4-bit as its authors describe it, although the code computes it in BF16.
For V4.1-Flash at 64K and above, re-scoring layers are counted over the candidate pool, as in its report \cite{DeepSeekAI2026V41Flash}; the reference code scores every entry, which adds 58.8\% to the MACs per token at 1M.

\subsection{Model Fidelity and the Source of Performance Data}\label{sec:fidelity}

%% P:sec5.5-p1
\ptag{sec5.5-p1}%
Single-PE latency comes from \evRTL; the fused placement and everything beyond one PE need a model \TBD{P}{source of performance data}.
Whichever model is used, its results are labeled as such and the workload inventory is an input, not a result.
A model calibrated against RTL is calibrated on \cfgAsBuilt, then on \cfgCorrected with new held-out points; reproducing a known fault of \cfgAsBuilt is agreement with the RTL, not a performance sample.

\subsection{Physical Implementation}\label{sec:physical}

%% P:sec5.6-p1
\ptag{sec5.6-p1}%
OpenROAD-flow-scripts, an open-source flow, implements the RTL on the ASAP7 predictive 7-nm PDK \citeTBD{ASAP7 PDK; open-source implementation flow}.
Area and timing come from synthesis and place-and-route runs under the stated library, constraint and SRAM-model assumptions; each is a \evPostRoute, not a silicon measurement.
Runs use a single process corner, and the liberty views of the SRAM macros are generated: their zero leakage and zero minimum period are gaps in the model, not zero cost.
Power is not reported until an activity-based method exists.
